\documentclass[11pt]{article}
\usepackage[T1]{fontenc}
\usepackage[utf8]{inputenc}
\usepackage{lmodern}
\usepackage[margin=1in]{geometry}
\usepackage{amsmath,amssymb,amsthm,booktabs}
\usepackage{microtype}
\usepackage{xurl}
\usepackage[colorlinks=true,linkcolor=black,citecolor=blue,urlcolor=blue]{hyperref}

\newtheorem{theorem}{Theorem}
\newtheorem{lemma}[theorem]{Lemma}
\newtheorem{corollary}[theorem]{Corollary}

\title{Decimal Divisor Digits and Prime Repunits: A Classification of OEIS A209933}
\author{Carlo Corti\\Independent researcher\\
\href{mailto:carlo.corti@outlook.com}{\texttt{carlo.corti@outlook.com}}}
\date{November 2026}

\begin{document}
\maketitle

\begin{abstract}
OEIS A209933 consists of the positive integers divisible by every digit
of every positive divisor. We classify these integers as thirteen small
values and the products of a decimal prime repunit by one of eleven fixed
coefficients. The proof reduces the divisor condition to quotient-digit
constraints indexed by the greatest common divisor with \(2520\).
Finite certificates and a transition argument cover arbitrary decimal
lengths. The classification and the identity obtained by deleting \(48\)
to give OEIS A182183 are verified in Lean. A finite enumeration independent
of the classification and the repunit catalog finds exactly \(46\) solutions
below \(10^{24}\). It proves that the term after \(264\) is the \(19\)-digit
repunit and extends the historical \(24\)-term prefix by \(22\) indexed
values. Conditional on published repunit primality proofs and reported
search coverage, the classification also determines the values and ranks
through index \(101\).
\end{abstract}

\noindent\textit{2020 Mathematics Subject Classification.}
Primary 11A63. Secondary 11A51, 11B83.

\noindent\textit{Key words and phrases.}
Decimal digits, divisibility, prime repunits, finite-state verification,
integer sequences, Lean.

\section{Introduction}\label{sec:intro}

For a positive integer \(d\), let \(D(d)\) be the set of digits of its
ordinary decimal expansion. Define \(P(n)\) for \(n\ge 0\) by requiring
\(n>0\) and
\[
 \text{for every positive }d\mid n\text{ and every }c\in D(d),
 \qquad c>0\ \text{and}\ c\mid n.
\]
The increasing enumeration of \(\{n:P(n)\}\), indexed from \(1\), is
\((a(k))\), OEIS A209933 \cite{oeis209933}. For example, \(48\) satisfies
\(P\). The union of the digit sets of its divisors is
\(\{1,2,3,4,6,8\}\), and these digits all divide \(48\). The digit \(6\)
in its divisor \(16\) need only divide the parent \(48\).

Write \(R_t=(10^t-1)/9\) for the decimal repunit of length \(t\ge 1\).
The A209933 entry at revision 27 lists \(24\) terms ending at \(264\),
conjectures \(a(25)=R_{19}\), and records that A209933 contains all
prime repunits, the values of A004022 \cite{oeis209933,oeis004022}.
On 29 September 2024, Brown reported that any further term exceeds
\(10^{11}\). His linked program uses a factor sieve followed by digit
checks on the divisors \cite{oeis209933,brown}. In A182183,
\(R_{19},2R_{19},\ldots,8R_{19}\) already occupied ranks \(24\)--\(31\),
credited to Charles R Greathouse~IV in 2012 \cite{oeis182183}.
These values give known members of A209933. Their ranks in A209933
require excluding intervening values under its weaker divisor-digit
condition. Section~\ref{sec:class} establishes the full classification,
and Section~\ref{sec:finite} proves the ranks in the adopted prefix
independently of that classification.

A209933 is a subsequence of A034838, whose membership test uses the
digits of the integer itself \cite{oeis209933,oeis034838}. Branicky's
automata for the corresponding counts by length, A346267, use
remainders and information about permitted digits \cite{oeis346267}.
The states used here instead record enough information to check the
digits emitted by several long divisions at once.

\section{Classification}\label{sec:class}

Let \(B=\{1,2,3,4,5,6,7,8,9,12,24\}\).

\begin{theorem}\label{thm:class}
For every nonnegative integer \(n\),
\[
 P(n)
 \quad\Longleftrightarrow\quad
 n\in B\cup\{15,48\}
 \ \text{or}\
 n=mR_t\text{ for some }m\in B\text{ and prime }R_t.
\]
\end{theorem}

\subsection{Quotient digits and a finite transition table}

Since \(M=\operatorname{lcm}(1,\ldots,9)=2520\), a nonzero digit divides
\(n\) exactly when it divides \(g=\gcd(n,M)\). If \(10\mid g\), then
\(10\mid n\), and the last digit of the divisor \(n\) is zero. Hence
a solution requires \(10\nmid g\). The \(48\) positive divisors of
\(2520=2^3 3^2 5\cdot7\) include \(18\) multiples of \(10\), leaving
\(30\) eligible classes. For an eligible class \(g\), put
\[
 S_g=\{c:1\le c\le9,\ c\mid g\},\qquad
 \mathcal C_g=\{x>0:g\mid x,\ D(x/d)\subseteq S_g\text{ for every }d\mid g\}.
\]
Here and below divisors of a positive integer are positive. For a true
term \(n\) and its actual class \(g\), every \(n/d\) with \(d\mid g\)
is a divisor of \(n\), so \(n\in\mathcal C_g\). The relaxed set
\(\mathcal C_g\) also admits \(x\) with \(\gcd(x,M)\ne g\).

Write \(u\) for a positive decimal prefix and append a digit \(c\).
For \(d\mid g\), long division gives
\begin{equation}\label{eq:division}
 \left\lfloor\frac{10u+c}{d}\right\rfloor
 =10\left\lfloor\frac{u}{d}\right\rfloor
   +\left\lfloor\frac{10(u\bmod d)+c}{d}\right\rfloor .
\end{equation}
The last floor lies between \(0\) and \(9\). A positive quotient of a
prefix is a prefix of the corresponding final quotient. In detail, if
\(u=\lfloor x/10^k\rfloor\), then
\(\lfloor u/d\rfloor=\lfloor \lfloor x/d\rfloor/10^k\rfloor\).
A zero quotient at this stage represents an empty prefix. All quotients
are positive once \(u\ge1000\). The bound \(d\le g\le504<1000\)
holds because a class without a factor \(5\) divides \(504\), while
a class with a factor \(5\) is odd and divides \(315\).

For \(u\ge1000\), use the state \(r=u\bmod g\). Because \(d\mid g\),
the new digit in \eqref{eq:division} is
\(\lfloor(10(r\bmod d)+c)/d\rfloor\). Thus a transition with
\(1\le c\le9\) is allowed exactly when every emitted digit lies in \(S_g\). In particular,
\(d=1\) emits \(c\) itself, forcing \(c\in S_g\). Its destination is
\((10r+c)\bmod g\), and the accepting state is \(r=0\).
For example, take \(g=d=12\). The prefix \(1\) has quotient zero,
meaning that no quotient digit has yet appeared. At prefix \(1333\),
the quotient is \(111\) and the remainder is \(1\). Appending \(2\)
emits the digit \(1\) in this division, giving quotient \(1111\) and
remainder \(0\) at \(13332\). The initialization tests all four-digit
prefixes by their integer quotients, including those with nonzero
initial remainder.

\begin{table}[htbp]
\centering
\small
\begin{tabular}{@{}lll@{}}
\toprule
Class \(g\) & Four-digit prefixes (remainder) & All allowed transitions\\
\midrule
\(1,\ldots,8\) & \(1111\cdot g\ (0)\) & \(0\xrightarrow{g}0\)\\
\(9\) & \(1199\ (2),\ 9999\ (0)\) & \(2\xrightarrow{9}2,\ 0\xrightarrow{9}0\)\\
\(12\) & \(1332\ (0),\ 1333\ (1)\) & \(1\xrightarrow{2}0,\ 1\xrightarrow{3}1\)\\
\(18\) & \(1999\ (1)\) & \(1\xrightarrow{9}1\)\\
\(24\) & \(2664\ (0),\ 2666\ (2)\) & \(2\xrightarrow{4}0,\ 2\xrightarrow{6}2\)\\
\(36\) & \(3999\ (3)\) & \(3\xrightarrow{9}3\)\\
Other eligible \(g\) & none & no reachable states\\
\bottomrule
\end{tabular}
\caption{Initialization and transitions for the quotient-digit test.
The row \(1,\ldots,8\) applies separately to each class. The parenthesized value is the
state \(u\bmod g\). An arrow's label is the appended digit. A state
with no arrow has no allowed extension.}
\label{tab:states}
\end{table}

Here is the complete prescription for checking the finite certificates.
For each eligible \(g\), form its divisors and \(S_g\). For each
\(u=1000,\ldots,9999\), include \(u\) in the initialization exactly if
every digit of \(\lfloor u/d\rfloor\) belongs to \(S_g\) for all
\(d\mid g\). This requires \(30\cdot9000\) prefix tests. For each
resulting remainder, try the nine digits in the transition rule.
Table~\ref{tab:states} records all \(16\) initialized prefixes and
all \(16\) allowed edges across the classes. Each destination occurs
among the initialized remainders of its class, which closes the table
under further extensions. For \(1\le u<1000\), test only multiples
of \(g\) with the same whole-quotient digit rule. The small accepted
values are \(g,11g,111g\) for \(g=1,\ldots,9\), \(12,132\) for
\(g=12\), \(15\) for \(g=15\), and \(24,48,264\) for \(g=24\). All
other classes have empty small lists. The project diagnostic
\verb|python3 -B check/language.py| checks these lists and the
transition table by direct quotient tests \cite{project}. For
example, \(999\in\mathcal C_9\), but its divisor \(27\) contains \(2\),
which does not divide \(999\). The full predicate can thus reject a
word that passes the quotient filter.

\begin{lemma}\label{lem:language}
For \(g\in B\), the set \(\mathcal C_g\) consists of \(gR_t\) for
\(t\ge1\), together with \(48\) when \(g=24\).
For \(g=15\), it consists of \(15\) alone.
For all remaining eligible classes it is empty.
\end{lemma}

\begin{proof}
At the four-digit start \(u\), each quotient is positive and, by the
truncation identity, its digits are a prefix of the corresponding
final quotient. Thus a word passing the quotient test starts in
Table~\ref{tab:states}. The induction invariant is that the state
equals the current prefix modulo \(g\) and that all quotient digits
read so far belong to \(S_g\). By \eqref{eq:division}, this condition
extends by one digit exactly along a listed edge. Every destination
again occurs among the initialized states, so induction covers any
number of appended digits in both directions. At termination, remainder
zero enforces \(g\mid x\).

For \(g=1,\ldots,8\) all accepting words repeat \(g\). For \(g=9\) the
loop at \(0\) gives the words of repeated \(9\)s, while the loop at
\(2\) never accepts. The \(g=18,36\) loops also never accept. For
\(g=12\), an accepted word of length \(h\ge4\) has a leading \(1\),
\(h-2\) copies of \(3\), and a final \(2\). At length four this is
\(1332\). Longer words start with \(1333\), repeat the loop on \(3\)
as needed and exit on \(2\). Increasing the length by one inserts a
\(3\) before the last digit, taking the value \(x\) to \(10x+12\).
Thus the value is \(12R_{h-1}\). For \(g=24\), the length-four word
is \(2664\). Longer words start with \(2666\), repeat the loop on
\(6\) and exit on \(4\), giving \(24R_{h-1}\) by \(x\mapsto10x+24\).
In these two classes the accepting state is terminal. The separate
small lists complete the claim.
\end{proof}

\subsection{Excluding composite repunits and proving the converse}

\begin{proof}[Proof of Theorem~\ref{thm:class}]
Suppose \(P(n)\). Lemma~\ref{lem:language} with \(g=\gcd(n,2520)\)
gives \(n=15\), \(n=48\), or \(n=gR_t\) with \(g\in B\). For \(t=1\)
the last form gives \(n=g\). If \(t\ge2\) and \(R_t=uv\) with
\(u,v>1\), then \(u,v\) are odd. For each \(d\mid g\), both \(gu/d\)
and \(gv/d\) divide the original \(n\). Their digits therefore lie
in \(S_g\), so \(gu,gv\in\mathcal C_g\). The exception \(48\) in
\(\mathcal C_{24}\) would give a factor \(2\) of the odd \(R_t\).
Hence \(u\) and \(v\) are both repunits of length at least two.
Their product ends in \(21\), since \(11^2\equiv21\pmod{100}\).
But \(R_t\) ends in \(11\), a contradiction. The argument uses only
divisors of \(n\) and the permitted digits for \(n\).

For the converse, \(B\) is closed under taking positive divisors.
Every member of \(B\) is divisible by its own nonzero digits: this is
immediate for the one-digit members and follows by inspection for
\(12\) and \(24\). If \(d\mid m\) with \(m\in B\), then \(d\in B\)
and each digit of \(d\) divides \(d\) and hence \(m\). Thus \(P(m)\)
holds for \(m\in B\). The divisors of \(15\) are \(1,3,5,15\), and their
digits divide \(15\). The digit check for \(48\) was given in
Section~\ref{sec:intro}.
Now let \(p=R_t\) be prime and \(m\in B\). Then \(p\ge11\), while the
prime factors of \(m\) belong to \(\{2,3,5,7\}\), so the divisors of
\(mp\) are exactly \(d\) and \(dp\) for \(d\mid m\). For \(d\le9\),
\(dR_t\) repeats the digit \(d\). For \(d=12\) or \(24\), its digits
are among \(\{1,2,3\}\) or \(\{2,4,6\}\), respectively. This follows
from \(12R_2=132\), \(24R_2=264\), and the patterns in the
proof of the lemma for
\(t\ge3\). In all cases each such digit divides \(d\).
The digits of \(d\) itself also divide \(d\). This proves \(P(mp)\).
\end{proof}

\section{The exact relation to A182183}

A182183 consists of positive integers \(n\) such that every positive
divisor \(d\) of \(n\) is divisible by every decimal digit of every
positive divisor \(e\) of \(d\) \cite{oeis182183}. A zero digit makes
the test fail, since \(0\nmid d\) for \(d>0\). In our notation the
condition is \(P(d)\) for every \(d\mid n\). One can also test
whether each divisor of \(n\) is divisible by its own digits.
Indeed, for \(e\mid d\mid n\), each digit of \(e\) then divides \(e\)
and hence \(d\). In the reverse direction use \(e=d\) in \(P(d)\).

\begin{corollary}\label{cor:relation}
For every nonnegative integer \(n\),
\[
 \bigl(\forall d\ge0,\ d\mid n\ \Longrightarrow\ P(d)\bigr)
 \quad\Longleftrightarrow\quad
 P(n)\ \land\ n\ne48 .
\]
Thus the value set of A182183 is obtained from A209933 by deleting \(48\).
\end{corollary}

\begin{proof}
The left-hand condition implies \(P(n)\) by taking \(d=n\), so it
also excludes \(n=0\). It fails at \(48\) since \(16\mid48\) and the
digit \(6\) of \(16\) does not divide \(16\). In the reverse
direction apply Theorem~\ref{thm:class}. The divisors of \(m\in B\)
stay in \(B\), and the proper divisors of \(15\) are in \(B\).
Every divisor of \(mR_t\) with \(R_t\) prime is either in \(B\) or has
the same prime-repunit form with coefficient in \(B\). Each of these
divisors satisfies \(P\).
\end{proof}

Since \(a(17)=48\), the term of A182183 at rank \(k\) equals \(a(k)\)
for \(1\le k<17\), and \(a(k+1)\) for every \(k\ge17\) for which
\(a(k+1)\) exists. The A182183 entry records the inclusion, the
exclusion of \(48\), and closure under divisors. Greathouse also
identifies its divisibility-minimal members greater than \(1\) with
A116692, the primes with one distinct decimal digit
\cite{oeis182183,oeis116692}. Corollary~\ref{cor:relation} identifies
the entire set difference. The earlier A182183 ranks cited in
Section~\ref{sec:intro} can therefore be transferred with the
one-position shift.

\section{Formal verification}

The Lean project \cite{project} proves Theorem~\ref{thm:class} as
\texttt{A209933.classification}. The definition \texttt{Predicate n}
requires \(0<n\) and, for every \(d\mid n\), tests \(0<c\) and \(c\mid n\)
for every \(c\) in \texttt{Nat.digits 10 d}. Positivity of \(n\) rules
out the divisor \(0\). Lean stores decimal digits in reverse order,
which preserves this universal test. The list \texttt{coefficients}
is \(B\). The recursive \texttt{repunit} starts at \(R_0=0\) and
satisfies \(R_{t+1}=10R_t+1\).
The theorem \texttt{A209933.repunit\_closed} gives the closed formula.
In the later module \texttt{Relations.lean}, \texttt{DivisorClosed n}
requires \texttt{Predicate d} for each \(d\mid n\) in the naturals.
The theorem \texttt{A209933.divisor\_closed\_iff} identifies it with
\(\texttt{Predicate n}\land n\ne48\), as in Corollary~\ref{cor:relation}.

For \(g>0\), \texttt{quotients\_iff} identifies \texttt{Quotients g x}
with the quotient-digit condition for all positive \(d\mid g\),
including divisors greater than \(9\).
The initial and small-number certificates use the Boolean test
\texttt{fastQ}. The implication \texttt{fastQ\_of\_quotients} checks
that every word satisfying the quotient condition also satisfies that
test. The Python diagnostic \path{check/language.py} checks
Table~\ref{tab:states} with the whole-quotient test described in
Section~\ref{sec:class}.
For the formal all-length implication, the invariant in
\texttt{quotients\_invariant} says that a surviving word has shape
\(gR_t+e\) with an allowed offset \(e\), or has one of the dead
remainders. The theorems \texttt{live\_step\_certificate} and
\texttt{dead\_step\_certificate} check closure of those cases.
The zero remainder test extracts the candidate shapes, and the
formal proof excludes composite repunits and proves the converse.

The builds used Lean 4.36.0 \cite{lean} and Mathlib at
\path{fc4f59e577feb1a4956671bc04ac18aaaf8cd842}, with transitive
dependencies pinned in the Lake manifest. The named theorems and
their transitive dependencies were exported and replayed in an
empty environment using Lean's own kernel. The retained build logs
report \texttt{propext} and \texttt{Quot.sound} for the repunit formula.
The classification and deletion equivalences also use
\texttt{Classical.choice}. Finite decisions use kernel reduction.
The audited proof dependencies consist of those standard axioms and
proved declarations. The attempted Lean4Lean check stopped at a
format incompatibility.

The formal scope covers the three named results. The proof of the rank
formula is given in Section~\ref{sec:catalog}. Assigning ranks to
particular large prime repunits additionally requires external catalog
coverage. The Lean theorems use \texttt{Nat.Prime} as a mathematical
predicate. Particular primality certificates and the numerical tests
have separate evidence. The full pinned source tree and the original
build, axiom and replay records are retained with the project. Rebuilding
requires acquiring the pinned dependencies and regenerating or retrieving
their build artifacts. A reviewer reported that compiling the short-case
certificate exhausted a container's memory limit of about \(3\) GB.
With the proof split by class, the reviewer reported a successful build
and replay using the pinned toolchain and Lean's own kernel. The statement
of the certificate was preserved.

\section{A complete finite check below \texorpdfstring{\(10^{24}\)}{10 to the 24th power}}\label{sec:finite}

\begin{theorem}\label{thm:finite}
Exactly \(46\) positive integers below \(10^{24}\) satisfy \(P\).
The first \(24\) agree with the historical A209933 prefix.
The remaining \(22\), in increasing order, are the products of \(R_{19}\)
by the increasing members of \(B\), followed by the products of
\(R_{23}\) by those same coefficients.
\end{theorem}

\subsection{Local primality certificates}

For \(p=R_{19}\) and \(p=R_{23}\), the audit used a full-order
Lucas criterion (see Lehmer \cite{lehmer}). For each \(p\), the
factors in the following decompositions were proved prime by trial
division through their integer square roots:
\begin{align*}
 R_{19}-1
   &=2\cdot3^2\cdot5\cdot7\cdot11\cdot13\cdot19\cdot37
     \cdot52579\cdot333667,\\
 R_{23}-1
   &=2\cdot5\cdot11^2\cdot23\cdot4093\cdot8779
     \cdot21649\cdot513239 .
\end{align*}
For \(p=R_{19}\) take the base \(b=3\), and for \(p=R_{23}\) take
\(b=11\). In both cases modular exponentiation gives
\[
 b^{p-1}\equiv1\pmod p,\qquad
 \gcd\bigl(b^{(p-1)/q}-1,p\bigr)=1
       \quad\text{for every prime }q\mid p-1.
\]
For a prime divisor \(\ell\) of \(p\), the first congruence gives
\(b^{p-1}\equiv1\pmod\ell\). Thus \(b\) is a unit modulo \(\ell\)
and its order divides \(p-1\). If this order were a proper divisor,
it would divide \((p-1)/q\) for some prime \(q\mid p-1\), which
contradicts the corresponding gcd condition. The order therefore
equals \(p-1\), and Lagrange's theorem gives \(p-1\mid\ell-1\).
Since \(2\le\ell\le p\), we get \(\ell=p\), proving primality.

\subsection{Exhaustive filtering and divisor checks}

The independent Python audit uses the elementary necessary
quotient-prefix condition but does not use the global classification,
the transition table, the repunit catalog or either materialization
program. For each eligible \(g\) it starts with the empty prefix and
appends digits from \(S_g\) up to length \(24\). It retains a prefix
\(u\) only if for every \(d\mid g\) the whole integer quotient
\(\lfloor u/d\rfloor\) is either zero or has all digits in \(S_g\).
For any solution of length at most \(24\) every one of its prefixes
survives in its true class, by the truncation identity. At each
length a surviving positive prefix with \(\gcd(u,2520)=g\) is also
tested as a complete candidate.

For each candidate \(n\), a rejected case has a witnessed divisor
with a zero digit or a digit failing to divide \(n\). The verifier
first tries \(21649\mid R_{11}\) and \(2071723\mid R_{17}\).
These factors exceed the trial-divisor ceiling of \(10000\).
Each is used only after testing that it divides the candidate and
contains an offending digit. In the recorded run they supplied
\(22\) and \(11\) rejections, respectively. The verifier then
tests divisor pairs exposed by trial divisors up to
\(\min(10000,\lfloor\sqrt n\rfloor)\). If that covers the entire
range through \(\lfloor\sqrt n\rfloor\) and no witness occurs,
the verifier certifies \(P(n)\). For each candidate still unresolved,
it checks for a split \(n=mp\) with one of the two locally certified
primes \(p\), \(m\le2520\), and \(\gcd(m,p)=1\). If the split succeeds,
it checks every divisor \(d\) and \(dp\) with \(d\mid m\).
The run aborts if the split or the complete acceptance check fails.
The bound \(10^{24}\) is the range over which this candidate
disposition was completed.

The execution attempted \(2456\) digit extensions and found \(182\)
distinct complete candidates. It recorded explicit invalidating
divisors for \(136\), and complete divisor checks succeeded for the
other \(46\). The sorted result agrees at every index with the
adopted b-file and at its first \(24\) indices with the separately
preserved historical table. As an additional check, we evaluated the two
lists \((mR_{19}:m\in B)\) and \((mR_{23}:m\in B)\), with \(m\) in
increasing order, by exact integer multiplication. Their concatenation
equals the last \(22\) entries of the adopted b-file.
Table~\ref{tab:ends} gives the endpoints. The completeness of the prefix
filter and the exhaustive disposition of the candidates
prove Theorem~\ref{thm:finite}, including the gap from \(264\) to
\(R_{19}\).

\begin{table}[htbp]
\centering
\begin{tabular}{@{}rll@{}}
\toprule
Index \(k\)&Form & \(a(k)\)\\
\midrule
24 & \(24R_2\) & \(264\)\\
25 & \(R_{19}\) & \(1111111111111111111\)\\
35 & \(24R_{19}\) & \(26666666666666666664\)\\
36 & \(R_{23}\) & \(11111111111111111111111\)\\
46 & \(24R_{23}\) & \(266666666666666666666664\)\\
\bottomrule
\end{tabular}
\caption{The historical endpoint and the two independently checked
blocks of the adopted prefix. Each block contains the eleven products
in coefficient order.}
\label{tab:ends}
\end{table}

To run the read-only audit, place the adopted b-file at
\path{data/b209933.txt} and the original historical table at
\path{outputs/b24.txt}, then run \verb|python3 -B check/audit.py|
from the project root. The standard-library program compares the
enumerated values with both tables after enumeration. Its JSON
output reports both Lucas certificates, all \(136\) divisor witnesses
and the sorted list of \(46\) solutions. The execution stops if a
frontier exceeds \(100{,}000\) prefixes or if an elapsed-time check
reaches the \(60\) second limit.

\section{Ordered blocks and the catalog derivation}\label{sec:catalog}

\subsection{Ordering and exact ranks under catalog coverage}

Let \(t_j\) be the \(j\)-th exponent for which \(R_{t_j}\) is prime,
where this rank exists, and write \(b_1<\cdots<b_{11}\) for the
elements of \(B\). If \(t=ab\) with \(a,b>1\), then
\(R_a>1\) is a proper divisor of \(R_t\) by the geometric-series
factorization. So every prime-repunit exponent is prime. The first
one is \(2\). The next is at least \(5\), since \(R_3=111\) is
composite. All exponents after \(2\) are odd, and adjacent prime
exponents there differ by at least two. Since
\(R_{t+2}=100R_t+11>24R_t\), blocks with successive prime-repunit
exponents are strictly separated.

The thirteen elements of \(B\cup\{15,48\}\) and the eleven
\(R_2\)-products are disjoint because the latter are multiples of
\(11\) and the former are not. Together they give exactly the first
\(24\) terms, ending at \(264\). Consequently, whenever \(j\ge2\)
and \(t_j\) exists,
\begin{equation}\label{eq:index}
 a(11j+2+i)=b_iR_{t_j}\qquad (1\le i\le11).
\end{equation}
Equation~\eqref{eq:index} gives the rank of a block once all earlier
prime repunits are known. Membership of a single block follows
from its primality and Theorem~\ref{thm:class}.

The longer derivation uses published primality proofs and reported
exclusion coverage. The base-10 Rprime status table for prime
exponents from \(3\) through \(10^6\) reports seven proved prime
repunits and one probable prime, at exponent \(270343\). It labels
the remaining cases factored or composite and reports zero unchecked
prime exponents \cite{rprime}. This supplies the external exclusion
premise for Table~\ref{tab:blocks}. The catalog input to the
materializer preserves A004023 as its original coverage source.
Batalov reported a wider search limit of \(10800000\) in that entry
\cite{oeis004023}. Panzone summarizes
the first five prime repunits \cite{panzone}. The proofs for
\(R_{317}\) and \(R_{1031}\) are due to Williams and to Williams
and Dubner, respectively \cite{williams317,williams}. Underwood's
\(R_{49081}\) proof is recorded by Primo \cite{primo}.
Enge certified \(R_{86453}\), and the later certification of
\(R_{109297}\) was carried out by Underwood and Enge
\cite{enge,enge109297,cm}. These proof records and the linked
certificate supplements are external to the local checks.

\begin{table}[htbp]
\centering
\begin{tabular}{@{}rll@{}}
\toprule
Exponent \(t\) & A209933 indices & Evidence used here\\
\midrule
\(2\)      & within \(1\)--\(24\) & \(R_2=11\), small case\\
\(19\)     & \(25\)--\(35\) & Local finite check and Lucas proof\\
\(23\)     & \(36\)--\(46\) & Local finite check and Lucas proof\\
\(317\)    & \(47\)--\(57\) & Published catalog and proof\\
\(1031\)   & \(58\)--\(68\) & Published catalog and proof\\
\(49081\)  & \(69\)--\(79\) & Published catalog and proof\\
\(86453\)  & \(80\)--\(90\) & Published catalog and proof\\
\(109297\) & \(91\)--\(101\) & Published catalog and proof\\
\bottomrule
\end{tabular}
\caption{Ordering of prime-repunit products. The first \(R_2\) block
interleaves with the small values. Ranks after \(46\) use the cited
primality and catalog premises.}
\label{tab:blocks}
\end{table}

Under these premises the last value of the auxiliary table is
\(24R_{109297}\) at rank \(101\). Under the same reported exclusion
coverage, a primality proof for \(R_{270343}\) would place its
products at ranks \(102\)--\(112\). In the November 2026 source
assessment \(R_{270343}\) remained a probable prime
\cite{oeis004023,rprime}. Infinitude of decimal prime repunits
remains open \cite{panzone}.

\subsection{Materialization and local validation}

The materializer emits one row at a time. It uses fixed decimal
patterns: a repeated digit for \(b_i\le9\), or the patterns for \(12\)
and \(24\) proved above. The checker constructs each product
digit by digit from the units place, adding the incoming carry to the
coefficient, taking the units digit, and passing the tens part to
the next iteration. Both routines share the admitted set of exponents
and the order of the blocks. This comparison checks arithmetic and
decimal serialization with work linear in the number of output
digits for fixed coefficients.

The 101-row output contains \(2{,}708{,}884\) bytes. In one recorded run,
generation and a subsequent standalone verification took about \(350\) ms
in total. The run receipt gives CPython 3.12.3, Linux kernel
6.17.0-42 and glibc 2.39. The two processes ran in sequence on the
author's workstation, identified by the author as an AMD Ryzen 9
7940HS running Linux Mint 22.3. The measurement covers decimal generation,
verification and file handling. The primality proofs, search for
prime-repunit exponents and formal verification have separate costs.

The 101-row table was generated and checked before the abrupt jump
after \(264\) prompted the separate finite audit in
Section~\ref{sec:finite}. In that later work, the full table was also
read into exact integers in nine-digit chunks. The audit checked
\(9n=m(10^t-1)\) for the repunit products after the historical prefix
and inspected \(643\) divisor occurrences across the \(101\) rows.
The large cases use the published primality inputs to justify complete
divisor lists. This audit supplies exact integer and membership
checks under those premises. Ranks after \(46\) continue to depend
on the reported catalog coverage. The author adopted the first \(46\)
terms as the canonical local result and the extent intended for OEIS.
The longer table documents the catalog-based derivation. The Lean
proof of the deletion identity in Corollary~\ref{cor:relation} was
completed after the scope decision.

\section{Conclusion}

The quotient-digit conditions force every solution into a small
exceptional set or a fixed multiple of a prime repunit. This rigidity
gives the deletion identity for A182183 and converts ranking after
the historical prefix into the ordering of prime-repunit exponents.
The finite audit establishes the conjectured term \(a(25)=R_{19}\)
and the exact adopted prefix through \(a(46)=24R_{23}\). A further
five blocks are accessible through published prime-repunit evidence
and its stated coverage.

\section*{Data availability}

The local publication package contains the adopted 46-term b-file,
the historical 24-term table, the quotient-table diagnostic, the finite
audit and its output, the catalog input, materializer and checker,
and essential execution and verification records. The pinned Lean
sources are supplied in a compact archive with a checked restoration
map; dependencies and build artifacts are acquired separately at the
recorded revisions. The package README gives fresh-root reproduction
commands and distinguishes retained historical proof checks from a
new source build. The auxiliary 101-row table is regenerated from
the supplied sources rather than duplicated in the package; its
historical receipt and exact-integer audit source and result are included.
Ranks beyond 46 retain the external premises stated above.
The GitHub URL in the references identifies the intended public
repository. This local preparation does not establish remote availability.
The provisional Zenodo DOI will be replaced with the verified
designated locator after the first Zenodo publication.

\section*{AI use disclosure}

Large language models were used as research and drafting assistants
during the A209933 workflow. OpenAI Codex assisted throughout the
project, and a Google Gemini preliminary analysis suggested the
quotient-digit filter. The separate Python code review and manuscript
review used external LLM-based review systems. Their reports identify
the ChatGPT, Claude, Gemini, Grok, Z.ai GLM, Meta and MiniMax
families across these activities. Exact model versions were not
consistently available in the review records.

\paragraph{Mathematical analysis.}
Codex assisted in developing the classification and its Lean proofs
and in assessing the source claims. The mathematical evidence is
the stated argument, the finite certificates and the audited Lean
theorems with pinned dependencies and replay logs.

\paragraph{Algorithm and code.}
Codex assisted with the Python materializer, the digit-by-digit
checker and the finite audit. The separate model-assisted code review
examined the Python artifacts, including their input handling and
verification behavior. The numerical claims rely on the
b-file, executable validation artifacts, Lucas certificates and
retained execution records. The project notes document conversational
assistance as part of the provenance.

\paragraph{Manuscript.}
Codex assisted with drafting, source reconciliation and internal checks.
Two external rounds of model-based manuscript review examined
the proof, reproducibility and attribution. These reviews were advisory.
The formal and computational claims are supported by the proofs,
programs and execution records described above.
The author chose the adopted numerical scope and takes full scientific
responsibility for the mathematical statements, algorithm,
implementation, numerical results, validation boundaries and
conclusions.

\clearpage
\begin{thebibliography}{99}
\raggedright\sloppy

\bibitem{oeis209933}
J. Krizek and contributors, \emph{OEIS A209933}, The On-Line
Encyclopedia of Integer Sequences, submitted 16 April 2012, revision 27
(29 September 2024). \url{https://oeis.org/A209933}.

\bibitem{brown}
L. A. Brown, \emph{A209933.py}, Python program, 2024.
\url{https://github.com/lucasaugustus/oeis/blob/main/A209933.py}.

\bibitem{oeis004022}
N. J. A. Sloane and contributors, \emph{OEIS A004022}, prime repunits,
revision 100 (2026). \url{https://oeis.org/A004022/internal}.

\bibitem{oeis182183}
J. Krizek and contributors, \emph{OEIS A182183}, The On-Line
Encyclopedia of Integer Sequences, submitted 17 April 2012, terms
\(24\)--\(31\) contributed by C. R. Greathouse IV on 17 April 2012.
\url{https://oeis.org/A182183}.

\bibitem{oeis034838}
E. Friedman and contributors, \emph{OEIS A034838}, integers divisible
by their own digits. \url{https://oeis.org/A034838/internal}.

\bibitem{oeis346267}
M. Marcus and contributors, \emph{OEIS A346267}, submitted 12 July
2021, with automaton constructions by M. S. Branicky.
\url{https://oeis.org/A346267/internal}.

\bibitem{project}
C. Corti, \emph{Decimal Divisor Digits and Prime Repunits: A
Classification of OEIS A209933}, GitHub repository and Zenodo archive,
2026. \url{https://github.com/carcorti/A209933},
\url{https://doi.org/10.5281/zenodo.xxxxxxxx}.

\bibitem{oeis116692}
R. L. Shepherd and contributors, \emph{OEIS A116692}, primes with
one distinct decimal digit, submitted 22 February 2006.
\url{https://oeis.org/A116692/internal}.

\bibitem{lean}
Lean FRO, \emph{The Lean Language Reference}, version 4.36.0.
\url{https://lean-lang.org/doc/reference/4.36.0/}.

\bibitem{lehmer}
D. H. Lehmer, \emph{Tests for primality by the converse of Fermat's
theorem}, Bull. Amer. Math. Soc. \textbf{33} (1927), no.~3, 327--340.

\bibitem{oeis004023}
N. J. A. Sloane and contributors, \emph{OEIS A004023}, exponents of
prime repunits, revision 299 (2026), with the search-limit report
by S. Batalov (2021). \url{https://oeis.org/A004023/internal}.

\bibitem{rprime}
K. Beschorner and Rprime project contributors, \emph{Rprime project
and base-10 exponent status table for \(3\) through \(10^6\)},
consulted 23 November 2026.
\url{https://www.kurtbeschorner.de/},
\url{https://www.kurtbeschorner.de/db-status-3-1M.htm}.

\bibitem{williams317}
H. C. Williams, \emph{Some primes with interesting digit patterns},
Math. Comp. \textbf{32} (1978), no.~144, 1306--1310.

\bibitem{williams}
H. C. Williams and H. Dubner, \emph{The primality of \(R_{1031}\)},
Math. Comp. \textbf{47} (1986), no.~176, 703--711.
\url{https://doi.org/10.2307/2008184}.

\bibitem{panzone}
P. A. Panzone, \emph{A remark on prime repunits},
Rev. Un. Mat. Argentina \textbf{53} (2012), no.~2, 31--35.
\url{https://inmabb.criba.edu.ar/revuma/pdf/v53n2/v53n2a04.pdf}.

\bibitem{primo}
M. Martin, \emph{Primo records}, entry for P. Underwood's
\(R_{49081}\) certificate (2022), page updated 25 February 2023.
\url{https://www.ellipsa.eu/public/primo/records.html}.

\bibitem{enge}
A. Enge, \emph{FastECPP over MPI}, in \emph{Mathematical Software:
ICMS 2024}, Lecture Notes in Comput. Sci. \textbf{14749},
Springer, 2024, 36--45.
\url{https://doi.org/10.1007/978-3-031-64529-7_4}.
Preprint (v2) \url{https://arxiv.org/abs/2404.05506v2}.

\bibitem{enge109297}
A. Enge, \emph{Primality record with ECPP -- 109297 digits},
author's report on joint work with P. Underwood (record set in 2025),
updated 11 February 2026.
\url{https://enge.math.u-bordeaux.fr/blog/ecpp-109297.html}.

\bibitem{cm}
A. Enge, \emph{ECPP}, CM project documentation and primality
certificate supplements, consulted November 2026.
\url{https://www.multiprecision.org/cm/ecpp.html}.

\end{thebibliography}
\end{document}
